\section{Material e methods}

\subsection{Study site and sampling strategy}

\begin{figure}[t]
    \begin{center}
    \includegraphics[width=\columnwidth]{LaTeX_Template_SBSR_English/figures/main_map_rework.png}
    \caption{a) Study site location with MapBiomas 2025 LULC map and sample polygons. Basemap: Google Satellite.}
    \label{fig:main_map}
    \end{center}
\end{figure}


The study site is the Brasilia National Park, located in the north-western region of the Federal District of Brazil (Fig.\ref{fig:main_map}).
The park is permanently protected and is a relevant research and conservation site \cite{sano2005}.
Its vegetation cover is representative of three major Cerrado physiognomies, namely Forest, Savanna, and Grasslands (Fig. \ref{fig:main_map}).
These physiognomies are referenced in this study as classes, and their spatial extents were extracted from the MapBiomas Collection 11 Land Use and Land Cover map of 2025 \cite{mapbiomas2026}.

We constructed a set of samples for each of the three Cerrado physiognomies.
The sample sets are composed of 5 sample polygons labeled after a color: White (W), Blue (U), Black (B), Red (R), and Green (G).
On a first approach of this study, we verified that a large number of observations ($n > 8000$ per sample) systematically leads to the rejection of the null hypothesis (H$_{0}$), whereas empirical evidence suggests otherwise.
Consequently, the decision was made to reduce the number of observations per sample to mitigate the effect of large sample sizes on the test statistic's p-value.
Hence, the sample polygons of each class were plotted as rectangles of \qtyproduct{150 x 200}{\meter} with 300 observations per polygon.
The polygons were distributed over the study site using the MapBiomas 2025 LULC map as reference (Fig. \ref{fig:main_map}), and were used to extract the observations ($O_{i}$) of intensity SAR data.
For Sentinel-1 images, $O_{i_{S1}} = (I_{VV}, I_{VH})$, and for NISAR images, $O_{i_{NISAR}} = (I_{HH}, I_{HV})$.

\subsection{SAR datasets and preprocessing}

We studied two scenes from two different SAR missions: a L-band NISAR Geocoded Single Look Complex (GSLC) image acquired in 2026-06-24 (Fig.\ref{fig:main_map}c), and a C-band Sentinel-1 Single Look Complex (SLC) image acquired in 2026-06-25 (Fig. \ref{fig:main_map}d).
The products of both mission were downloaded from the Alaska Satellite Facility (ASF) platform.

For both missions, we used the intensity data represented by the diagonal elements of the C2 matrix ($C_{11}$ and $C_{22})$ expressed in power linear units.
The Sentinel-1 image contains the VV and VH intensity polarizations and was submitted to TOPSAR Split, orbit file correction, TOPSAR deburst, multilook processing with a 4x1 window, C2 matrix generation and terrain correction with pixel resampling to 10m.
The NISAR image contains de HH and HV intensity polarizations, and was submitted to a multilook processing with a 2x2 window resulting on a \qtyproduct{10 x 10}{\meter} resolution image, followed by the C2 matrix processing.

\subsection{Experimental Design}

Under a statistical analysis approach, a series of parametric hypothesis tests was employed to  verify whether samples were produced by the same model for both Sentinel-1 and NISAR intensity data over Cerrado physiognomies.
The study unfolds in five sequential steps, where one step leads to the next one: 1. Estimate the equivalent number of looks (ENL) of both NISAR and Sentinel-1 scenes; 2. Goodness-of-fit tests for two models: $\Gamma_{SAR}$ and $\mathcal G^0_I$; 3. Apply a stochastic distances hypothesis test to verify the separability of the samples; 4. For those samples in which the null hypothesis was rejected, we further verified if this was due to a difference in the means only, or on the roughness only, or on both parameters; 5. Given the main difference between the samples, we further verified if fixing a parameter and freeing the other produces the same distribution.

For all hypothesis tests, the adopted significance level was $1\%$.
This research was executed in \url{Quarto} version 1.10.18 and is fully reproducible.
All data, code, and experimental results are available on \url{github.com/JalesBussinguer/sar_image_analysis_sbsr}.

To estimate the equivalent number of looks (ENL) of the scenes, we plotted square polygons of \qtyproduct{200 x 200}{\meter} distributed all over water bodies and exposed soil regions around the scenes.
The ENL was obtained as the reciprocal of the square of the sample variation coefficient.
For NISAR 40 polygons were used, and the ENL was 3.3795 for C11 and 3.4487 for C22.
For Sentinel-1, 20 polygons were used, and the ENL was 2.0574 for C11 and 2.0788 for C22.
Those values are used as parameters for the tested models and are fixed.

To assess which model describes better the samples intensity distributions, we proceeded to goodness-of-fit tests for the distributions $\Gamma_{SAR}$ and $\mathcal G^0_I$.
The $\Gamma_{SAR}$ model is a reparameterization of the classical gamma law that characterizes fully developed speckle \cite[Eq. 1]{alpala2025}.
The $\mathcal G^0_I$ distribution is well recognized for its capacity to describe different levels of heterogeneity \cite[Eq. 5]{alpala2025}.

To test the goodness-of-fit for both distributions, we used a parametric bootstrap over the Kolmogorov-Smirnov (KS) test statistic $D$.
Since the parameters are estimated from the sample itself, the null distribution of $D$ is not the one tabulated for KS; therefore, the quantiles are obtained via simulation (parametric bootstrap) with 1000 repetitions, re-estimating the parameters using maximum likelihood in each replicate.
After the tests, we adopted the $\mathcal G^0_I$ as the model that better characterizes the sample distributions.
The decision was based on evidences about the adherence to data and its tractability and interpretability of its parameters.

We then proceeded to the comparison of the same-class samples' distributions, to verify their similarity/separability.
Given the absence of tests to verify whether two samples originate from the same distribution with an exact distribution under the H$_{0}$, we used the symmetrized Kullback-Leibler (KL) stochastic distance \cite[item i.]{nascimento2010} hypothesis test defined by \cite[Eq. 9]{nascimento2010}: $S_{\mathrm{KL}}$.
The test's H$_{0}$ is the equality between the parameters of the sample distributions.
This family of hypothesis tests presents well suited specificity and sensitivity [citação?].

The $S_{\mathrm{KL}}$ test only indicates whether two sample distributions differ, but it does not qualify the rejection of H$_{0}$.
To achieve this qualification, we followed the methodology described by \cite{nascimento2010}, evaluating the pairs of classes' distributions across four distinct scenarios based on their shape ($\alpha$) and scale ($\mu$) parameters:
(i) the distributions share the same shape parameter but differ in their scale parameter ($\alpha_1=\alpha_2$ and $\mu_1\neq\mu_2$);
(ii) the scale parameters are identical, whereas the shape parameters vary ($\alpha_1\neq\alpha_2$ and $\mu_1=\mu_2$);
(iii) both the shape and scale parameters of the first distribution are strictly smaller than those of the second ($\alpha_1<\alpha_2$ and $\mu_1<\mu_2$);
(iv) the first distribution exhibits a smaller shape parameter but a larger scale parameter compared to the second ($\alpha_1<\alpha_2$ and $\mu_1>\mu_2$).
The pairs of sample distributions which the H$_{0}$ was rejected were separated, being the focus of this study from this point.

We used two likelihood ratio tests with 1 degree of freedom: one for $\alpha_1=\alpha_2$ (and $\mu_1,\mu_2$ free), and another for $\mu_1=\mu_2$ (with $\alpha_1,\alpha_2$ free).
When neither tests rejects H$_{0}$, but $S_{\mathrm{KL}}$ rejects, the difference is classified as combined effect.
With this framework it is possible to understand the reason behind the H$_{0}$ rejection in the $S_{\mathrm{KL}}$ test.

Based on the results of the previous tests, which points that the differences are mainly related to the mean brightness, the following question arise: it is possible that the sample distributions share the same roughness parameter, where only the mean brightness varies?
This question seek evidence that supports the idea of a class being characterized by a family of distributions that share the same roughness parameter.
To do so, we used a likelihood ratio test for each class and polarization.
With 5 samples per class, the null hypothesis is H$_{0}:\alpha_1=\dots=\alpha_5$, with $\mu_1,\dots,\mu_5$ free, against the alternative that all $\alpha_j$ and $\mu_j$ are free.
The test statistic $-2\log\Lambda$ have asymptotic distribution $\chi^2_{k-1}$, with 4 degrees of freedom.
Rejecting H$_{0}$ at the significance level of $1\%$ supports the idea of a family of distributions that characterize the class under analysis.